\bookchapter{第五十五章：專案二——天氣查詢 App（網路版）}

\chapterintro{記帳本是“練內功”，天氣 App 是“見世面”——它要和網際網路打交道。這個專案你會學到：申請免費 API、解析 JSON、處理“載入中/成功/失敗”三種介面狀態，以及怎麼把 API 金鑰藏好。做完它，你就具備了做任何“聯網 App”的能力——新聞、匯率、翻譯，套路一模一樣。}

\booksection{一、申請免費天氣 API}

個人練手，推薦 OpenWeatherMap 或國內的 和風天氣（QWeather）——都有免費額度，註冊即用。以 OpenWeatherMap 為例：

\begin{enumerate}
  \item 註冊帳號 在 API keys 頁面生成一個 Key（apikey）；
  \item 免費版可呼叫 Current Weather（當前天氣）介面：api.openweathermap.org/data/2.5/weather? q=beijing\&appid=你的Key\&units=metric；
  \item 瀏覽器裡先試一次，看看返回的 JSON 長什麼樣——這是接任何 API 前必做的"預習"。
\end{enumerate}

\textbf{避坑提醒：金鑰絕不能寫死在程式碼裡}\par

API Key 相當於你的“錢包密碼”，寫進程式碼一提交到 GitHub，幾小時內就會被爬蟲掃走盜刷。正確做法：放進 local.properties（這個檔案已被 gitignore 忽略），再透過 Gradle 的 buildConfig 讀進程式碼。搜尋“Android 隱藏 API Key 的最佳實踐”照著改即可。

\booksection{二、網路層：Retrofit + JSON 解析}

先照 API 返回的 JSON 定義“資料模型”，欄位對不上的一律忽略：

\begin{lstlisting}[language=Kotlin]
// JSON: {"weather":[{"description":"晴"}],"main":{"temp":26.5,"humidity":40},"name":"北京"}
data class WeatherResponse(
    val weather: List<Weather>,
    val main: Main,
    val name: String
)
data class Weather(val description: String)
data class Main(val temp: Double, val humidity: Int)

interface WeatherApi {
    @GET("data/2.5/weather")
    suspend fun getWeather(
        @Query("q") city: String,
        @Query("appid") key: String,
        @Query("units") metric: String = "metric",
        @Query("lang") lang: String = "zh_cn"
    ): WeatherResponse
}
\end{lstlisting}

Gson 會自動把 JSON 塞進這些 data class——欄位名一致，魔法自動發生。

\booksection{三、三種介面狀態：載入中 / 成功 / 失敗}

聯網 App 的靈魂在於“狀態管理”。網路慢、斷網、城市名打錯……每種情況介面都要給使用者交代。用一個 sealed class（密封類）優雅地表示：

\begin{lstlisting}[language=Kotlin]
sealed class UiState {
    object Loading : UiState()                    // 轉圈中
    data class Success(val data: WeatherResponse) : UiState()  // 成功
    data class Error(val message: String) : UiState()          // 失敗
}

// ViewModel 裡：
fun search(city: String) {
    viewModelScope.launch {
        _state.value = UiState.Loading
        _state.value = try {
            UiState.Success(api.getWeather(city))
        } catch (e: Exception) {
            UiState.Error("查不到這個城市，換個名字試試？")
        }
    }
}

// 介面裡：
when (val s = state) {
    is UiState.Loading -> CircularProgressIndicator()   // 轉圈
    is UiState.Success -> WeatherContent(s.data)        // 展示
    is UiState.Error -> Text(s.message)                 // 報錯
}
\end{lstlisting}

這個 Loading/Success/Error 三段式是聯網介面的“萬金油模板”，先背下來，以後所有聯網 App 照葫蘆畫瓢。

\booksection{四、實戰介面程式碼：把天氣畫出來}

狀態有了，介面長這樣（關鍵片段）：

\begin{lstlisting}[language=Kotlin]
@Composable
fun WeatherContent(data: WeatherResponse) {
    Column(
        Modifier.fillMaxSize().padding(24.dp),
        horizontalAlignment = Alignment.CenterHorizontally
    ) {
        Text(data.name, fontSize = 28.sp, fontWeight = FontWeight.Bold)
        Text(
            "${data.main.temp}°C",
            fontSize = 72.sp,
            fontWeight = FontWeight.Light
        )
        Text(
            data.weather.firstOrNull()?.description ?: "未知",
            fontSize = 20.sp,
            color = Color.Gray
        )
        Spacer(Modifier.height(16.dp))
        Row {
            Text(" 溼度 ${data.main.humidity}%")
            Spacer(Modifier.width(24.dp))
            Text(" 體感 ${data.main.temp}°C")
        }
    }
}
\end{lstlisting}

幾個小細節值得學：firstOrNull() 安全取第一個元素（列表可能為空），?: “未知” 給兜底文案——空安全原則在介面裡同樣適用。溫度用 72.sp 的大字號撐場面，“輕”字重（Light）讓數字顯得優雅，這是天氣類 App 的經典配方。

\booksection{五、介面與互動細節}

\booksection{七、介面與互動細節}

\begin{itemize}
  \item 搜尋框：頂部 TextField，輸城市名點搜尋（鍵盤上設"搜尋"動作，回車即搜）；
  \item 天氣卡片：大字號溫度（26.5°）、天氣描述（晴）、溼度、城市名；
  \item 收藏城市：把常看的城市存進 DataStore，下次開啟自動查——順手複習第 6 章；
  \item 下拉重新整理：Compose 裡用 PullToRefreshBox（Material 3 提供），手感立刻"專業"起來。
\end{itemize}

\begin{tcbinfo}{小貼士：圖片讓介面活起來}
給每種天氣配一個 emoji 或圖示（），用 when(description) 對映——三行程式碼，介面溫度感翻倍。也可以載入第三方的天氣圖示庫，但練手階段 emoji 足夠。

\end{tcbinfo}

\begin{tcbpractice}{來動手：天氣 App 驗收清單}
(1) 正常搜尋出天氣；(2) 斷開網路，App 優雅報錯不崩潰（開飛航模式測！）；(3) 輸一個不存在的城市，提示友好；(4) 收藏兩個城市，重啟 App 還在；(5) 轉屏不丟狀態。五條全過，這個專案才算畢業。

\end{tcbpractice}

\begin{tcbsummary}{本章小結}
\begin{itemize}
  \item 先瀏覽器試 API、看懂 JSON，再寫程式碼——"預習"永遠不虧
  \item API Key 藏進 local.properties，絕不進程式碼庫
  \item Retrofit 介面 + data class：欄位對上，JSON 自動解析
  \item Loading/Success/Error 三段式狀態是聯網介面萬金油
  \item 斷網、空輸入、城市不存在——錯誤處理是聯網 App 的及格線
\end{itemize}

\end{tcbsummary}
